\section{Въведение: Епохата след персоналните компютри и новите изчислителни изисквания}

Процесът на трансформация на формите на човешкия труд поставя качествено нови проблеми за решаване от едно общество.
Компютърните технологии като незаменим инструмент в решаването на по-голямата част от тези проблеми, също трябва
да се адаптират подходящо, за да откликнат на нуждите на новите стопански и социални реалности.
Нещо повече, оказва се, че промените в компютърните технологии, подобно на социално-икономическите изменения следват
принципите на еволюцията и ко-еволюцията.

И наистина, първите разработени компютри са програмирани на машинен код,
по-късно, за улеснение машинните инструкции били заменени с мнемонични, които са тяхно биективно съответствие.
Машинният код и мнемоничните езици (асемблер) като семантика на инструкциите са много по-близо до програмирането на
физическата памет и хардуера на един компютър и в този смисъл те не са удобни за програмирането на абстрактни процеси.
Това неудобство е породило нуждата да се изгради една по-абстрактна архитектура от инструкции (\acr{ISA} - Instruction Set Architecture), така, че
да се описват процеси, които да са от по-широк научен интерес.

Тъй като първите научни кръгове, в които компютърните
технологии са станали достояние са били тези на математиците и инженерите, то естествено е да се очаква, че първите
езици за програмиране след мнемоничните ще имат \acr{ISA}, което е изключително удобно за разработката и програмирането на
алгоритми и математически изчисления. Така първите езици от висок ред като \acr{Fortran} (FORmula TRANslation, 1957) и \acr{Algol} (ALGOrithmic Language, 1958) били създадени предимно
за изграждането на алгоритми и математически изчисления. Те използвали близък до математическия език \acr{ISA}, който
посредством разбор се ``превеждал'' или както е утвърдено да се нарича ``компилирал'' до асемблер, а от там до машинен код.

Оказало се обаче, че тези езици са не дотам удобни при програмирането на задачи от стопански и производствен характер
(писане на операционна система, правене на програма за складови наличности, четене на данни от много входове и др.).
Търсенията на програмистите за удачен език след този период съвсем не следвали някакви закономерности, а по-скоро
имали случайностен характер. Например езикът Pascal бил разработен от Niklaus Wirth през 1970 г. с цел да служи за преподаване на студентите на структурно програмиране,
\acr{Ada} бил разработен от Министерството на отбраната на САЩ през 1980 г. за критични военни и авиационни системи с високи изисквания за надеждност и безопасност,
\acr{BASIC} (Beginner's All-purpose Symbolic Instruction Code) бил разработен от John Kemeny и Thomas Kurtz през 1964 г. в Dartmouth College като прост език за обучение на начинаещи програмисти.

Най-сполучлив се
оказал езикът за програмиране C, създаден от Dennis Ritchie в Bell Labs през 1972 г. Този език притежавал най-добрата
комбинация от \acr{ISA}, които са удобни и за алгоритмичен и за хардуерен дизайн. С него се пишело удобно както алгоритмично
поставен проблем, така и програмен проблем (като например писането на операционна система. Даже компилаторът на C е написан на самия език C - така наречения ``self-hosting compiler'').
Освен това, той бил и достатъчно ефективен откъм бързодействие. Тук пак компилаторът изиграл много важна роля,
защото той имал за задача да сведе инструкциите на C до инструкции на асемблер.

Благодарение на езика C се създали
много и различни по вид езици като например Prolog (1972, логическо програмиране), C++ (1985, обектно-ориентиран C), Perl (1987, текстова обработка), Python (1991, скриптове и обща употреба), Ruby (1995, уеб приложения), и много други. Някои от езиците приличали по синтаксис
и функционалности на C / C++, а други се стремели да се разграничат от неговата семантика.

Заслужава внимание езикът Java (създаден от James Gosling в Sun Microsystems през 1995 г.),
който е подобен на C++, но кодът се компилира не до машинен, а до байткод (bytecode), който след това се изпълнява от Java Virtual Machine (\acr{JVM}) за всяка една компютърна архитектура, така че да е максимално преносим (``write once, run anywhere''). Освен това при Java не се допуска директно менажиране на паметта (автоматичен garbage collector).

Други езици се отдалечили от парадигмата на C, като JavaScript (създаден от Brendan Eich в Netscape през 1995 г. за 10 дни) и Python например. Те напълно премахнали
строгото статично типизиране на данните и се съсредоточили върху \acr{ISA}, който да е оптимален за бързо прототипиране, динамична обработка на данни и интерактивна разработка.

Постепенно в практиката се наложили езиците, които не използват строги статични типове от данни и са с опростен синтаксис и не съдържат
или съдържат ограничен брой директни машинни инструкции (високо ниво на абстракция). Практиката и характерът на програмните продукти утвърдила парадигмата
на обектно-ориентираното програмиране пред останалите. Ако пътят на програмирането беше праволинеен
или подвластен на някакъв закон би се очаквало тези тенденции да се затвърждават с течението на времето.
Случва се обаче точно обратното.

След утвърждаването на уеб технологиите и мобилните платформи (епохата ``Post-\acr{PC}'' след 2010 г.), потребността от по-ефективно изпълнение на
даден процес отново връща на мода строго типизираните езици, като най-разпространени стават Java и C\#.
JavaScript откликва на това предизвикателство като се превръща и в backend език чрез технологията Node.js (създадена от Ryan Dahl през 2009 г.,
по-късно и чрез Deno и Bun) и чрез ``псевдотипизирането'' си в езика TypeScript (създаден от Microsoft през 2012 г., добавящ статични типове към JavaScript).

Python откликва чрез използването
на \acr{JIT} (Just-In-Time) компилатори като Numba или PyPy, както и чрез нови компилатори като Cython, като също въвежда опционално типизиране (type hints) от версия 3.5 (2015 г.).

Но историята съвсем не свършва до тук. Писането на все по-ефективни програми, обработващи големи данни (Big Data),
става все по-голямо предизвикателство и едно от най-простите работещи решения е паралелната обработка.
Отначало паралелната обработка се свеждала до нишково програмиране (multi-threading), впоследствие се разрастнало до
паралелно свързани изчислителни машини (cluster computing, distributed computing), а в днешно време се е утвърдила парадигмата на хетерогенното програмиране,
с която изпълнението на програмите се разпределя между централния процесор (\acr{CPU}) и графичния процесор (\acr{GPU}).

Това доведе отново до ``възцаряването'' на езика C (и C++), който поддържа програмирането на графична карта чрез технологии като \acr{CUDA} (NVIDIA, 2007), \acr{OpenCL} (Khronos Group, 2009), \acr{SYCL} и протокол за
управление на паметта на графичната карта, а строго типизираните данни са задължително условие
при изпълняването на код върху графичната карта (заради оптимизациите на хардуерно ниво и \acr{SIMD} - Single Instruction Multiple Data).

Езиците като Python намират решение на този проблем като
създават компилатор за хетерогенни програми с помощта на \acr{LLVM} (като Numba за \acr{CUDA} kernels), а други като TypeScript, чрез \acr{FFI} (Foreign Function Interface) протокол за извикване на C/C++ код.
За отбелязване е, че всички езици ползват \acr{FFI} протокол или подобни механизми, понеже само C/C++ (и други системни езици) могат да управляват директно паметта на графичната карта и да извикват \acr{CUDA}/\acr{OpenCL}/\acr{Vulkan} \acr{API}-та.

Внимателният читател без труд ще забележи, че и в този случай
развитието на езиците за програмиране не следва някакъв строго определен праволинеен ход, а е низ от случайности,
потребности на обществото и резерв на изменчивост, който езиците притежават, за да откликнат на нуждите на
нововъзникващата среда. Така ходът на развитието на програмните езици има по-скоро зигзагообразен вид и
е еволюционен по своята природа, а не детерминистичен.

Друго достойно за отбелязване явление е, че отново възниква
нуждата от компилатори и инфраструктура за компилация (като \acr{LLVM}), чрез които да се оптимизира ефективността на кода за различни архитектури (x86, \acr{ARM}, \acr{RISC-V}, \acr{GPU}).

Това разбиране за компютърните технологии и интерпретация на събитията в сферата на програмирането
аз до голяма степен получих от крупния научен работник, аристократ и общественик Петър Свещаров при многобройните беседи, които сне водили. 
Благодарение на него аз се вдъхнових и изградих този проект.

